\subsection{ms-contenidos --- Archivos (S3), módulos e ítems}
\label{sec:ms-contenidos}
\textit{Responsable: Alexander (CON-01, CON-02). Repositorio: \url{https://github.com/nubeduca/nubeduca-ms-contenidos}.}

\pendiente{Dos o tres párrafos: qué resuelve este microservicio, decisiones de diseño y qué roles (STUDENT, TEACHER, ADMIN) usan cada operación.}

\subsubsection{Endpoints}
\pendiente{Tabla con método, ruta, roles y evento que emite. Sale del OpenAPI del microservicio en la carpeta \texttt{apis/} de \texttt{nubeduca-contratos} (campos \texttt{x-roles} y \texttt{x-emits}).}

\subsubsection{Eventos}
\textbf{Publica:}
\begin{itemize}
  \item \texttt{lms.content.file.scanned.v1}
  \item \texttt{lms.content.module.published.v1}
\end{itemize}

\textbf{Consume:}
\begin{itemize}
  \item \texttt{lms.courses.course.created.v1}
  \item \texttt{lms.courses.enrollment.created.v1}
  \item \texttt{lms.courses.enrollment.dropped.v1}
\end{itemize}

\pendiente{Alexander: explicar el flujo de archivos (URL prefirmada, evento de S3, Lambda de antivirus, estados CLEAN/INFECTED) y el bucket usado.}

\subsubsection{Estructura JSON de la tabla \texttt{<stage>-t\_contenidos}}
Clave de partición \texttt{tenant\_id}, clave de ordenamiento \texttt{sk}.

\begin{lstlisting}
{
  "tenant_id": "utec",
  "sk": "<PREFIJO>#<id>",
  "...": "completar con un ítem real de la tabla"
}
\end{lstlisting}
\pendiente{Reemplazar por un ítem real de cada tipo que guarda la tabla (por ejemplo con \texttt{aws dynamodb scan --table-name dev-t\_contenidos --max-items 1}). Una tabla con los prefijos de \texttt{sk} y un bloque JSON por tipo de ítem.}

\subsubsection{Datos de prueba}
\pendiente{Evidencia de los 30\,000 registros (UTEC 15\,000, UPC 10\,000, PUCP 5\,000): salida de \texttt{poblar\_datos.py} o un conteo por tenant con \texttt{aws dynamodb query ... --select COUNT}.}

\subsubsection{Resultados y evidencias}
\evidencia{ms-contenidos-evidencia.png}{ms-contenidos: prueba funcional (por ejemplo, una llamada en Postman o una captura del frontend).}{fig:ms-contenidos}
